\documentclass[a4paper, serif, 10pt]{moderncv}

%% moderncv
% color and style
\moderncvcolor{black}
\moderncvstyle{banking}

%% page layout/margins
\usepackage[top=2.5cm, bottom=2.5cm, left=1.5cm, right=1.5cm]{geometry}

\nopagenumbers{}

%% Babel config for Brazilian Portuguese language
% ref:
% - https://latex3.github.io/babel/guides/locale-portuguese.html
\usepackage[brazilian]{babel}

%% fontspec
\usepackage{fontspec}

\IfFontExistsTF{Linux Libertine}{
  \setmainfont[Ligatures={TeX, Common}, Numbers=OldStyle]{Linux Libertine}
}{}
\IfFontExistsTF{Linux Libertine O}{
  \setmainfont[Ligatures={TeX, Common}, Numbers=OldStyle]{Linux Libertine O}
}{}

%% CTeX
% CJK support, used to show CJK characters in the resume
%
% - heading=false: disable ctex heading
% - punct=kaiming: use kaiming punctuations styles for CJK
% - scheme=plain: use plain scheme, do not override \normalsize font size
% - space=auto: space settings for CJK characters
%
% CJK fonts are resolved per locale language, e.g. Noto Serif/Sans CJK JP
% for Japanese resumes, with fallback to the next region if not installed
%
% ref:
% - http://ctan.mirrorcatalogs.com/language/chinese/ctex/ctex.pdf
\usepackage[UTF8, heading=false, punct=kaiming, scheme=plain, space=auto]{ctex}

\IfFontExistsTF{Noto Serif CJK SC}{
  \setCJKmainfont{Noto Serif CJK SC}
}{}
\IfFontExistsTF{Noto Sans CJK SC}{
  \setCJKsansfont{Noto Sans CJK SC}
}{}

%% line spacing
\usepackage{setspace}
\setstretch{1.125}

%% URL styling - use normal text instead of monospace
\urlstyle{same}

%% Patch \httplink and \httpslink to handle full URLs
% The original moderncv macros always prepend http:// or https:// to URLs.
% This causes issues when the URL already has a protocol (e.g., https://example.com)
% resulting in malformed URLs like https://https://example.com.
% This patch checks if the URL already contains "://" and uses it directly if so.
% Note: We use \str_set:Nx (x-expansion) to expand macros like \@homepage first.
\ExplSyntaxOn
\str_new:N \g_yamlresume_url_str

\RenewDocumentCommand{\httpslink}{O{}m}{
  \str_set:Nx \g_yamlresume_url_str {#2}
  \str_if_in:NnTF \g_yamlresume_url_str {://}
    { \url{#2} }
    { \url{https://#2} }
}

\RenewDocumentCommand{\httplink}{O{}m}{
  \str_set:Nx \g_yamlresume_url_str {#2}
  \str_if_in:NnTF \g_yamlresume_url_str {://}
    { \url{#2} }
    { \url{http://#2} }
}
\ExplSyntaxOff

\name{Mariana Costa Almeida}{}
\title{Engenheira de Dados Sênior | Plataformas de Dados e Analytics}
\phone[mobile]{+55 11 98765-4321}
\email{mariana.costa.almeida@example.com}
\homepage{https://marianaalmeida.dev}

\address{Rua dos Pinheiros, 450, Apto 82, São Paulo, SP, Brasil, 05422-000}{}{}

\extrainfo{{\small \faLinkedin}\ \href{https://www.linkedin.com/in/mariana-costa-almeida}{@mariana-costa-almeida} {} {} {} • {} {} {} 
{\small \faGithub}\ \href{https://github.com/mariana-costa-data}{@mariana-costa-data}}

\begin{document}

\maketitle

\section{Informações Básicas}

\cvline{}{\begin{itemize}
\item Engenheira de dados com mais de 6 anos de experiência na construção de plataformas de dados escaláveis, confiáveis e orientadas a impacto de negócio
\item Especialista em pipelines batch e streaming, modelagem analítica, qualidade de dados e governança em ambientes de nuvem
\item Experiência com Apache Spark, Airflow, dbt, Kafka, AWS, Databricks e ferramentas modernas de lakehouse
\item Atuação colaborativa com times de produto, ciência de dados, BI e engenharia para traduzir necessidades analíticas em soluções robustas
\item Forte compromisso com boas práticas de engenharia, observabilidade, automação e mentoria técnica
\end{itemize}}
\section{Formação}

\cventry{mar. de 2014–dez. de 2018}
        {Bacharelado, Ciência da Computação, Nota: 8.7/10}
        {Universidade de São Paulo}
        {\url{https://www.usp.br}}
        {}
        {\begin{itemize}
\item Formação sólida em fundamentos de computação, bancos de dados e engenharia de software
\item Desenvolvimento de projetos acadêmicos com foco em processamento e análise de grandes volumes de dados
\item Participação em grupos de estudo de ciência de dados e computação distribuída
\end{itemize}
\textbf{Cursos}: Estruturas de Dados, Algoritmos e Complexidade, Bancos de Dados, Engenharia de Software, Sistemas Distribuídos, Estatística Aplicada, Aprendizado de Máquina, Computação em Nuvem}

\cventry{mar. de 2019–dez. de 2020}
        {Mestrado, Ciência de Dados, Nota: 9.1/10}
        {Insper}
        {\url{https://www.insper.edu.br}}
        {}
        {\begin{itemize}
\item Especialização em arquitetura de dados, modelagem analítica e pipelines escaláveis
\item Projeto de conclusão sobre qualidade de dados em plataformas de analytics em tempo real
\item Desenvolvimento de soluções com foco em impacto de negócio e boas práticas de engenharia
\end{itemize}
\textbf{Cursos}: Engenharia de Dados, Estatística para Ciência de Dados, Aprendizado de Máquina, Visualização de Dados, Governança de Dados}
\section{Experiência Profissional}

\cventry{fev. de 2022–Atual}
        {Engenheira de Dados Sênior}
        {Nubank}
        {\url{https://nubank.com.br}}
        {}
        {\begin{itemize}
\item Liderança técnica na construção e evolução de pipelines de dados batch e streaming para produtos de crédito e risco
\item Modelagem de dados analíticos em lakehouse com Apache Iceberg, Spark e dbt, apoiando times de analytics e ciência de dados
\item Implementação de contratos de dados, testes automatizados e monitoramento de qualidade com foco em confiabilidade
\item Otimização de custos e performance em AWS, reduzindo tempo de processamento e consumo de recursos
\item Mentoria de engenheiras e engenheiros de dados júnior e colaboração com produto, engenharia e negócio
\end{itemize}
\textbf{Palavras-chave}: Apache Spark, Apache Airflow, dbt, Kafka, AWS, Data Quality, Lakehouse}

\cventry{jan. de 2020–jan. de 2022}
        {Engenheira de Dados}
        {iFood}
        {\url{https://www.ifood.com.br}}
        {}
        {\begin{itemize}
\item Desenvolvimento de pipelines ETL e ELT para ingestão de dados de pedidos, entregadores e restaurantes
\item Construção de modelos dimensionais e tabelas analíticas para dashboards de operações e growth
\item Automação de orquestração com Airflow e padronização de transformações com SQL e Python
\item Atuação em iniciativas de observabilidade de dados, catálogo e linhagem para maior confiabilidade
\item Parceria com times de BI e produto para traduzir necessidades analíticas em soluções escaláveis
\end{itemize}
\textbf{Palavras-chave}: ETL, ELT, SQL, Python, Airflow, Data Warehouse}

\cventry{mar. de 2019–dez. de 2019}
        {Analista de Dados Júnior}
        {Itaú Unibanco}
        {\url{https://www.itau.com.br}}
        {}
        {\begin{itemize}
\item Manutenção e evolução de rotinas de carga de dados para relatórios regulatórios e gerenciais
\item Criação de consultas SQL e scripts Python para tratamento, validação e consolidação de dados
\item Apoio à documentação de processos, regras de negócio e dicionários de dados
\item Colaboração com áreas de risco e finanças para garantir consistência das informações
\end{itemize}
\textbf{Palavras-chave}: SQL, Python, Dados Financeiros, Relatórios, Qualidade de Dados}
\section{Idiomas}

\cvline{Português}{Nativo ou Bilíngue \hfill \textbf{Palavras-chave}: Comunicação técnica}
\cvline{Inglês}{Proficiência Profissional Fluente \hfill \textbf{Palavras-chave}: TOEFL iBT 108, Inglês para negócios}
\cvline{Espanhol}{Proficiência Profissional Limitada \hfill \textbf{Palavras-chave}: Leitura técnica}
\section{Competências}

\cvline{Engenharia de Dados}{Especialista \hfill \textbf{Palavras-chave}: Apache Spark, Apache Airflow, dbt, Apache Kafka, Databricks, Apache Iceberg, Modelagem de Dados, Data Quality}
\cvline{Nuvem e Infraestrutura}{Avançado \hfill \textbf{Palavras-chave}: AWS, Terraform, Docker, Kubernetes, Amazon S3, AWS Glue, Snowflake}
\cvline{Programação e Analytics}{Avançado \hfill \textbf{Palavras-chave}: Python, SQL, Scala, PySpark, Git, Power BI, Looker}
\section{Prêmios}

\cventry{nov. de 2023}
        {Nubank}
        {Prêmio de Excelência em Engenharia de Dados}
        {}
        {}
        {Reconhecimento por liderança técnica na modernização da plataforma de dados de risco, com ganhos de confiabilidade e redução de custos operacionais.}
\section{Certificados}

\cventry{mar. de 2023}
        {Amazon Web Services}
        {AWS Certified Data Analytics - Specialty}
        {\url{https://aws.amazon.com/certification/certified-data-analytics-specialty/}}
        {}
        {}

\cventry{ago. de 2022}
        {Google Cloud}
        {Google Cloud Professional Data Engineer}
        {\url{https://cloud.google.com/learn/certification/data-engineer}}
        {}
        {}
\section{Publicações}

\cventry{ago. de 2023}
        {Construindo uma Plataforma de Dados Confiável em Escala}
        {Blog de Engenharia do Nubank}
        {\url{https://medium.com/@mariana.costa.almeida/plataforma-de-dados-confiavel}}
        {}
        {\begin{itemize}
\item Artigo sobre práticas de contratos de dados, testes automatizados e observabilidade em pipelines críticos
\item Discussão de arquitetura lakehouse, governança e métricas de confiabilidade para times de engenharia
\end{itemize}}
\section{Referências}

\cventry{\emaillink[ricardo.menezes@insper.edu.br]{ricardo.menezes@insper.edu.br}}
        {Professor de Ciência de Dados}
        {Prof. Ricardo Menezes}
        {+55 11 91234-5678}
        {}
        {Mariana combina rigor técnico, visão de negócio e capacidade de liderança. Seu trabalho em plataformas de dados trouxe ganhos concretos de confiabilidade e eficiência para equipes multidisciplinares.}
\section{Projetos}

\cventry{mai. de 2023–nov. de 2023}
        {Plataforma de streaming para ingestão, enriquecimento e scoring de eventos transacionais}
        {Pipeline em Tempo Real para Detecção de Fraudes}
        {\url{https://github.com/mariana-costa-data/fraud-streaming-pipeline}}
        {}
        {\begin{itemize}
\item Arquitetura com Kafka, Spark Structured Streaming e armazenamento em lakehouse para análise de eventos em near real time
\item Implementação de regras de qualidade, alertas e monitoramento de latência e throughput
\item Integração com modelos de machine learning para apoio à decisão em risco e prevenção a fraudes
\end{itemize}
\textbf{Palavras-chave}: Kafka, Spark Streaming, Lakehouse, Fraude, Data Quality}
\section{Interesses}

\cvline{Corrida}{Corrida de rua, Meia maratona, Trilhas}
\cvline{Comunidade}{Mentoria, Palestras, Open Source}
\cvline{Música}{Violão, MPB, Jazz}
\section{Voluntariado}

\cventry{mar. de 2021–Atual}
        {Mentora Voluntária}
        {Data Girls Brasil}
        {\url{https://www.datagirlsbrasil.org.br}}
        {}
        {\begin{itemize}
\item Mentoria de mulheres em transição de carreira para dados, com foco em SQL, Python e projetos práticos
\item Participação em eventos e rodas de conversa sobre carreira em engenharia de dados
\item Apoio na revisão de currículos e preparação para entrevistas técnicas
\end{itemize}}
    
\end{document}